\documentclass[11pt]{article}

\usepackage[margin=1in]{geometry}
\usepackage{amsmath}
\usepackage{booktabs}
\usepackage{enumitem}
\usepackage{graphicx}
\usepackage{siunitx}
\usepackage{hyperref}
\usepackage{gensymb}


\title{Intermediate Crawling Grade Test}
\author{Rift / Team 09}
\date{\today}

\begin{document}


\maketitle

\section*{Test ID}
\textbf{T-01ICb Intermediate Crawling Grade Test}

\section*{Test Objective}
Quantify the maximum incline the payload can successfully maneuver. Testing will consider both lateral and direct motion on inclined surfaces.

\section*{Robot Configuration}
\begin{itemize}[noitemsep]
    \item Robot configurations: Standing
    \item Tube pressures: 15 psi
    \item Control mode: Open-loop
    \item Electronics status: Batteries are fully charged
\end{itemize}

\section*{Materials}
\begin{itemize}[noitemsep]
    \item Stopwatch timer
    \item Measuring tape (at least 5 m in length)
    
\end{itemize}

\section*{Environment and Setup}
\begin{itemize}[noitemsep]
    \item Surface type: Concrete, Asphalt, Grass, etc.
    \item We will utilize different environments around campus, including the engineering loading platform, \underline{\hspace{5 cm}}
    \item Surface angle: $\psi \neq 0  {\degree}$ 
    \item External load: None
    \item Measurement method: time the rover moving directly and laterally through 2 full step cycles on different grades
\end{itemize}

\section*{Initial Conditions}
\begin{itemize}[noitemsep]
    \item Initial position: $x_0 = 0$
    \item Initial orientation: 
    \begin{itemize}
        \item Direct crawling: $\theta_z = 0 {\degree}$
        \item Lateral crawling:  $\theta_z = 90 {\degree}$
    \end{itemize}
    \item Initial shape state: Standing
    \item Stabilization time before command: 10?~s
\end{itemize}

\section*{Command Sequence}
\begin{itemize}[noitemsep]
    \item Command type: \underline{\hspace{5 cm}}
    \item Command values: \underline{\hspace{5 cm}}
    \item Update rate: 3~Hz
    \item Command Steps: 8~s
\end{itemize}

\section*{Measured Quantities}
\begin{itemize}[noitemsep]
    \item Speed of the rover on different inclines
    \item Grades provided from Facilities Management office
\end{itemize}

\section*{Success Criteria (Per Trial)}
\begin{itemize}[noitemsep]
    \item Forward traversal speed: .1 km/h  $\leq$ .14 km/h $\leq$ No upper limit
    \item Side to Side traversal speed: ?? km/h  $\leq$ ?? km/h $\leq$ No upper limit
    \item  Step height: 10 cm  $\leq$ 25 cm $\leq$ No upper limit
\end{itemize}

\section*{Test Procedure}
\begin{enumerate}[noitemsep]
    \item Place robot in standing position on a graded surface
    \item Complete pretest checklist
    \item     Send command sequence to drive rover 8 steps (2 full step cycles) forward while starting timer
    \item Stop timer upon completion of sequence
    \item Record what grade the rover was walking on
    \item Repeat steps 1-5, driving the rover laterally instead of directly on the same incline

\end{enumerate}

\section*{Number of Trials}
\begin{itemize}[noitemsep]
    \item Total trials: 1
    \item Time between trials: N/A
\end{itemize}

\section*{Results}
Summarize the observed performance metrics.

\begin{center}
\begin{tabular}{llcllll}\toprule
 Incline Angle (degrees)& \multicolumn{2}{c}{} & \multicolumn{2}{c}{}& \multicolumn{2}{c}{}\\\midrule
 & Direct Crawling ($\dot x$)&Lateral Crawling ($\dot y$) & Direct Crawling ($\dot x$)& Lateral Crawling ($\dot y$)& Direct Crawling ($\dot x$)&Lateral Crawling ($\dot y$)\\

Traversal Speed (km/h)&   & & & & &\\ \bottomrule
\end{tabular}
\end{center}

\section*{Observed Failure Modes}
Describe any observed failure behaviors.

\begin{itemize}[noitemsep]
    \item Truss buckles during stepping
    \item Truss fails to lift nodes from the ground
    \item Truss tips while lifting node from the ground
\end{itemize}

\section*{Conclusions}
Provide a concise interpretation of the results.

\textit{Example:} The robot achieved consistent forward motion on flat ground with an average speed of \underline{\hspace{ 1cm}}~cm/s. Performance degraded beyond a slope of \underline{\hspace{1 cm}}~\si{\degree}, where sustained forward progress was not achieved.

\section*{Notes and Limitations}
Document sources of variability or limitations.

\end{document}